\chapter{Root Canal Treatment}

\section*{Overview}
Root canal treatment, or endodontic therapy, removes infected or damaged pulp from inside a tooth, cleans and seals the root canals, and preserves the natural tooth. The exact approach, setting, and preparation depend on the indication, anatomy, and the patient's health.

\section*{Alternative Names}
Endodontic therapy; root canal procedure; endodontics

\section*{Principle}
After removing the pulp chamber and canals, the dentist shapes and disinfects the canal system and fills it with a biocompatible material (gutta-percha), preventing reinfection and allowing the tooth to be restored.

\section*{History}
Root canal therapy was developed in the late 19th and early 20th centuries with the use of gutta-percha and root canal instruments, becoming a routine dental procedure.

\section*{Why It Is Done}
Performed for irreversible pulpitis, pulp necrosis, dental abscess, and trauma or fracture exposing the pulp, to relieve pain and save the tooth.

\section*{Is This a Primary or Secondary Test or Procedure}
Primary definitive treatment for pulp and periapical disease.

\section*{How It Is Performed}
Performed under local anesthesia and rubber dam isolation. The tooth is opened, canals are shaped with rotary files, irrigated with disinfectant, and obturated with gutta-percha. A permanent crown or filling follows.

\section*{Preparation}
Radiographs and vitality testing; antibiotics are reserved for systemic signs of infection.

\section*{Contraindications and Precautions}
A non-restorable tooth, gross vertical root fracture, and inadequate bone support are contraindications.

\section*{Risks}
Postoperative pain, instrument fracture, incomplete obturation, reinfection, and tooth fracture without a crown.

\section*{Aftercare and Recovery}
Mild soreness is expected; the tooth is restored with a crown to prevent fracture.

\section*{Outcome and Follow-up}
Successful therapy yields a healed apex with no pain, swelling, or fistula at follow-up radiographs.

\section*{Expected Results}
A symptom-free tooth with intact root filling and no periapical radiolucency.

\section*{Follow-up}
Clinical and radiographic review; retreatment or extraction if symptoms or lesions persist.

\section*{When to Seek Medical Care}
Follow the procedure team's specific instructions. Seek urgent medical assessment for severe or worsening pain, uncontrolled bleeding, fever or signs of infection, trouble breathing, chest pain, fainting, new weakness or confusion, inability to urinate, or any rapidly worsening symptom. The appropriate warning signs depend on the procedure and the patient's medical condition.

\section*{References}
\begin{enumerate}
  \item MedlinePlus. Root canal. U.S. National Library of Medicine; 2025.
  \item Current Medical Diagnosis and Treatment. McGraw-Hill; 2025.
\end{enumerate}

\clearpage
